% Generated draft: requires mathematical and source review.
\documentclass[11pt]{article}
\usepackage[margin=25mm]{geometry}
\usepackage{fontspec}
\setmainfont{Latin Modern Roman}
\usepackage{amsmath,amssymb,mathtools,unicode-math}
\setmathfont{Latin Modern Math}
\usepackage{enumitem,needspace,xurl,hyperref}
\setlength{\parindent}{0pt}
\setlength{\parskip}{5pt}
\setlength{\emergencystretch}{3em}
\newcommand{\N}{\mathbb{N}}
\newcommand{\Z}{\mathbb{Z}}
\newcommand{\Q}{\mathbb{Q}}
\newcommand{\R}{\mathbb{R}}
\newcommand{\C}{\mathbb{C}}
\newcommand{\F}{\mathbb{F}}
\newcommand{\A}{\mathbb{A}}
\newcommand{\PP}{\mathbb{P}}
\newcommand{\E}{\mathbb{E}}
\newcommand{\eps}{\varepsilon}
\newcommand{\dd}{\,\mathrm{d}}
\begin{document}
\section*{UnsolvedMath ID: 30000438}
\textbf{Source problem code:} catalogue entry
\subsection{Definitions and mathematical statement}
The source asks whether the following are equivalent for a graph $G$:
\[
 \exists\text{ a curved-edge embedding }f:|G|\hookrightarrow\R^3
       \text{ with no strong quadrisecant line}
 \quad\Longleftrightarrow\quad G\text{ is }1/2\text{-planar}.
\]
Here $|G|$ denotes the geometric realization of the graph. An embedding requires edges to meet only at their prescribed common endpoints. Curved edges allow more than straight-line realizations, but the source does not specify smooth, piecewise-linear, tame or other regularity, nor explicitly the finite/simple graph convention.

An ordinary quadrisecant is an affine line meeting the embedded graph in four distinct points. The additional word \emph{strong} is not defined in the record: requirements on distinct edges, ordering, transverse intersection or separation cannot be guessed. Likewise $1/2$-planar is a source-specific class whose definition is absent. It is not silently identified with planar, $1$-planar, a prescribed crossing-number bound or a half-plane drawing condition. Consequently the displayed equivalence accurately retains the target but needs both definitions before it becomes an independently complete classification statement.

The literature triage checked 21 August 2026 says the quadrisecant characterization remains open in its converse direction; it does not reproduce the established implication or identify its precise hypotheses. The same assessment separately reports impossibility for a torus-triangulation question with one degree-$5$ and one degree-$7$ vertex, associated with Izmestiev--Rote--Sullivan. That separate result does not resolve either direction of this spatial-graph characterization, and its full triangulation hypotheses are not supplied in this entry. The registry partially-solved label and these two distinct status claims are preserved without conflating their objects.
\subsection{Short English statement}
Are the graphs admitting a curved spatial embedding without a strong quadrisecant exactly the source's $1/2$-planar graphs? Both specialized definitions are missing. The supplied assessment leaves a converse open and separately cites a solved torus-triangulation question, which is not a solution of this equivalence.
\subsection{Sources}
\begin{itemize}
\item Local UnsolvedMath registry, permanent ID 30000438, exact source code catalogue entry; \emph{Quadrisecant-Free Spatial Graph Embeddings}. Proposed year 2006, status partially\_solved; corrected\_verified statement retains the curved-edge, strong-quadrisecant and $1/2$-planarity terms. Verification's report check is metadata, not an inspection performed here.
\item \emph{Discrete Differential Geometry (2006)}, Oberwolfach Reports, DOI 10.4171/owr/2006/12, pp. 653--728, \url{https://doi.org/10.4171/owr/2006/12}.
\item Supplied literature triage checked 2026-08-21: ``original OWR problems'', \url{https://doi.org/10.4171/owr/2006/12}; ``Izmestiev--Rote--Sullivan'', \url{https://arxiv.org/abs/1207.3605}. The latter is cited in the assessment for the separate torus-triangulation impossibility statement. This dated OWR-derived curation requests independent verification of time-sensitive status and cited claims.
\end{itemize}
\end{document}
